\documentclass[11pt]{article}
\usepackage[margin=1in]{geometry}
\usepackage{hyperref}
\usepackage{listings}
\usepackage{xcolor}
\usepackage{enumitem}
\usepackage{booktabs}

\lstset{
  basicstyle=\ttfamily\small,
  breaklines=true,
  columns=fullflexible,
  frame=single,
  backgroundcolor=\color{gray!10},
  showstringspaces=false
}

\title{OPeNDAP and pydap: A Working Summary\\
       \large grep-dap study phase, Prompt \#1}
\author{Compiled for P. Cornillon}
\date{2026-05-16}

\begin{document}
\maketitle

\begin{abstract}
This document summarizes OPeNDAP (the Data Access Protocol, or DAP) and
the \texttt{pydap} Python client, with emphasis on the three aspects
that matter for inferring the contents of a remote OPeNDAP server:
\emph{metadata}, \emph{data retrieval}, and \emph{subsetting}. The
intent is to give \texttt{grep-dap} a precise, mechanical understanding
of what an OPeNDAP server exposes, how that exposure differs between
DAP2 and DAP4, and how \texttt{pydap} surfaces those exposures to a
Python program. Where the underlying documentation flags details as
out-of-date or fragile, the same caveats are repeated here.
\end{abstract}

\section{What OPeNDAP is, briefly}

OPeNDAP is a client--server protocol for remote access to scientific
datasets. A server publishes datasets at URLs; a client constructs a
URL that names a dataset and, optionally, a \emph{constraint
expression} that restricts what the server should compute and return.
The on-wire payloads are \emph{compressed}, the data are
\emph{remotely subset before transmission}, and the protocol enforces a
rigid syntactic description of dataset structure while leaving
semantic description (units, provenance, long names, conventions)
largely unconstrained. The asymmetry is intentional: it lowered the
barrier to publication at the cost of letting semantic metadata vary
from complete to absent across servers~--- which is the very problem
that motivates \texttt{grep-dap}.

Two versions of the protocol coexist in deployment:
\begin{description}[leftmargin=2em,style=nextline]
  \item[DAP2 (\textasciitilde 1993--2010s).] The original protocol.
    Carries a flat data model with arrays, \texttt{Grid}s (an array
    bundled with its coordinate axes), \texttt{Structure}s, and
    \texttt{Sequence}s. No groups. No 64-bit integers. Two metadata
    documents: \emph{DDS} (structure) and \emph{DAS} (attributes).
    Binary data extension is \texttt{.dods}.
  \item[DAP4 (2016--).] Jointly specified by OPeNDAP, Inc.\ and
    Unidata. Roughly a superset of DAP2 \emph{except} that
    \texttt{Grid} is replaced by an array with shared dimensions and
    \texttt{Map}s. Adds \texttt{Group}s (NetCDF-4 / HDF5--style
    hierarchy), fully qualified names (FQNs), 64-bit integers,
    \texttt{Opaque}, \texttt{Enum}, chunked HTTP/1.1 transfer, and
    checksums. One metadata document: \emph{DMR} (Dataset Metadata
    Response). Binary data extension is \texttt{.dap}.
\end{description}

A DAP4-capable server (Hyrax, modern THREDDS) generally also speaks
DAP2 by flattening or translating. The translation can fail: if a
dataset contains a type DAP2 cannot represent (e.g.\ \texttt{Int64}),
the DAP2 response is an HTTP 400 error explaining as much.

Two server implementations dominate:
\texttt{Hyrax} (OPeNDAP, Inc.) and \texttt{THREDDS} / \texttt{TDS}
(Unidata). Many NASA Earthdata holdings sit behind Hyrax; CMIP6 and
many academic/agency archives sit behind THREDDS. ERDDAP, while not a
DAP server in the strict sense, exposes a DAP-compatible endpoint
(\texttt{tabledap}, \texttt{griddap}) that \texttt{pydap} can read.

\section{Metadata}

For \texttt{grep-dap}, metadata is the first thing on the wire and the
cheapest thing to fetch. Understanding which document a server returns
and what it contains is the foundation of everything else.

\subsection{The shape of a DAP URL}

Every OPeNDAP dataset is addressable as a base URL. Appending a known
extension selects a \emph{response} from that dataset. The most useful
extensions, by protocol:

\begin{center}
\begin{tabular}{@{}lllp{0.45\linewidth}@{}}
\toprule
Ext. & Protocol & Content & Notes \\
\midrule
\texttt{.dds}      & DAP2 & Structure (text)              & ``Dataset Descriptor Structure''~--- shapes, types, hierarchy. \\
\texttt{.das}      & DAP2 & Attributes (text)             & ``Dataset Attribute Structure''~--- semantic metadata. \\
\texttt{.dods}     & DAP2 & Headers + binary              & Carries the actual data. \\
\texttt{.ascii} / \texttt{.asc} & DAP2 & ASCII data        & Human-readable data dump. \\
\texttt{.html}     & either & Data Request Form           & The browser form for building constraint expressions. \\
\texttt{.dmr}      & DAP4 & DMR (XML)                     & ``Dataset Metadata Response''~--- structure + attributes in one document. \\
\texttt{.dmr.xml}  & DAP4 & DMR (XML, explicit)           & Same content, explicit MIME. \\
\texttt{.dmr.html} & DAP4 & DMR rendered                  & Browser-friendly view of the DMR. \\
\texttt{.dap}      & DAP4 & Headers + binary (chunked)    & Carries the actual data; supports checksums. \\
\bottomrule
\end{tabular}
\end{center}

A reasonable rule of thumb when probing an unfamiliar URL:
\begin{enumerate}
  \item Append \texttt{.dmr} (or \texttt{.dmr.html}) and try. If you
    get back XML or a rendered page, the server speaks DAP4 and you
    should prefer it.
  \item If \texttt{.dmr} 404s, append \texttt{.dds} and \texttt{.das};
    if those return, the server is DAP2-only.
  \item For a THREDDS URL the path segment \texttt{dodsC} indicates a
    DAP2 endpoint; replacing it with \texttt{dap4} gives the DAP4
    endpoint when one exists.
\end{enumerate}

\subsection{DAP2 metadata: DDS + DAS}

DAP2 splits structure and attributes:

\textbf{DDS (Dataset Descriptor Structure)} is a small text document
that names variables, their data types
(\texttt{Byte}, \texttt{Int16}, \texttt{UInt16}, \texttt{Int32},
\texttt{UInt32}, \texttt{Float32}, \texttt{Float64}, \texttt{String},
\texttt{URL}), and their shape. Container types are
\texttt{Structure}, \texttt{Sequence}, and \texttt{Grid}. A
\texttt{Grid} bundles an N-D array with N 1-D ``map'' arrays giving
its coordinate axes; this lets a Grid be self-describing in the way
NetCDF intends.

\textbf{DAS (Dataset Attribute Structure)} is a parallel text document
listing attributes attached to each variable and to the dataset as a
whole. The DAS is where \texttt{units}, \texttt{long\_name},
\texttt{standard\_name}, \texttt{\_FillValue}, \texttt{valid\_range},
provenance strings, CF convention tags, and arbitrary free-text
metadata live. The DAS is rigid in syntax only; its keys and values
are entirely at the data provider's discretion. \emph{This is where
OPeNDAP's loose semantic constraints live.}

\subsection{DAP4 metadata: the DMR}

DAP4 collapses DDS and DAS into a single XML document, the DMR. The
DMR adds:
\begin{itemize}
  \item \texttt{Group}s, which may be nested. A dataset always has a
    root group. Each group can hold its own variables, dimensions,
    and attributes. Groups have a single parent (no cycles).
  \item Shared \emph{dimensions} declared once at the group level and
    referenced by variables. Variables list their dimensions, by
    name, in order.
  \item \emph{Fully qualified names} of the form
    \texttt{/group/subgroup/var} that uniquely identify a variable
    regardless of name collisions in sibling groups.
  \item Atomic types missing from DAP2:
    \texttt{Int64}, \texttt{UInt64}, \texttt{Enum}, \texttt{Opaque}.
  \item \texttt{Map}s: a variable can declare which other variables
    serve as its coordinate axes, recovering DAP2's
    \texttt{Grid}-style self-description without bundling.
\end{itemize}
Because the DMR carries both structure and attributes, a single
\texttt{HTTP GET} of \texttt{.dmr} gives a client everything it needs
to introspect a dataset without touching the data.

\subsection{Information available even when semantic metadata is poor}

Even when a DAS or DMR carries only a sparse set of attributes, an
inspecting client has the following to work with:
\begin{itemize}
  \item The variable \emph{names} themselves (\texttt{SST},
    \texttt{analysed\_sst}, \texttt{lat\_ph}, \texttt{u\_wind}, \ldots).
  \item Dimension names and ordering.
  \item Variable shapes; a 2-D field over (lat, lon) is shaped
    differently from a 1-D trajectory or a tabular sequence.
  \item Data types (a \texttt{Float32} 3-D field is structurally
    different from an \texttt{Int8} mask).
  \item Group structure and file/directory hierarchy.
  \item Whatever attributes are present (often at least
    \texttt{units}, \texttt{\_FillValue}, and perhaps a
    \texttt{long\_name} or \texttt{standard\_name}).
  \item The dataset's neighbors on the same server (via catalog
    pages), which give context.
  \item Small samples of the data, fetched via subsetting (next
    section), from which statistical and structural properties can
    be estimated cheaply.
\end{itemize}

\section{Data retrieval}

A DAP server returns data only in response to a request that names
the variables (and slices) wanted. The binary response is
\texttt{.dods} for DAP2 and \texttt{.dap} for DAP4. Each binary
response is preceded by an inline copy of the metadata describing
exactly what is in the payload (so a client never has to guess), then
the binary blob. DAP4 chunks the binary using HTTP/1.1 transfer
encoding and may include CRC-style checksums; DAP2 sends one
unchunked binary blob and has no checksum mechanism.

\subsection{The constraint expression}

A constraint expression (CE) is appended to the data URL with
\texttt{?} and tells the server what to compute. A request with no CE
asks for the entire dataset~--- almost always wrong for big files.

\textbf{DAP2 CE syntax} (variable projection plus optional sequence
selection):
\begin{lstlisting}
data_url?<var1>[i0:s:i1][j0:s:j1]...,<var2>...&<filter1>&<filter2>...
\end{lstlisting}
The square-bracketed triples are start, stride, end (\emph{inclusive}
in DAP2; this surprises people coming from Python). Filters apply to
\texttt{Sequence}s and look like \texttt{seq.profile\_id=5}.

\textbf{DAP4 CE syntax} (named parameter, FQNs, shared dimensions):
\begin{lstlisting}
data_url?dap4.ce=<FQN1>[i0:s:i1];<FQN2>[i0:s:i1];...
\end{lstlisting}
Note: \emph{semicolons separate variables in DAP4}, not commas.
DAP4 hyperslab indices are \texttt{[start:stride:stop]} but
\emph{stop is exclusive} in some server implementations and inclusive
in others~--- in practice, stop is treated as the last index to
include, matching DAP2. Always check with a tiny request.

\textbf{Shared-dimension CE in DAP4} (often what you actually want):
\begin{lstlisting}
data_url?dap4.ce=<FQN_dim1>=[i0:s:i1];<FQN_dim2>=[j0:s:j1];<FQN_var1>;<FQN_var2>;...
\end{lstlisting}
The dimension slices are defined once with \texttt{=}, and any listed
variable that uses those dimensions inherits the slice. This is far
less error-prone than restating the hyperslab on every variable.

\subsection{Server-side functions}

DAP servers may expose named functions that the client can invoke as
part of the CE; Hyrax, for example, exposes \texttt{geogrid} for
lat/lon bounding-box subsets in geographic coordinates rather than
index space. There is no standard discovery mechanism for these
functions; a client either knows the function exists (out-of-band) or
does not use it.

\subsection{What ``compressed and sent'' actually means}

The DAP2 \texttt{.dods} payload is the requested variables, with their
DDS prepended, serialized in XDR-like network byte order. The DAP4
\texttt{.dap} payload is similar but chunked, with the DMR for the
constrained dataset prepended and optional CRC checksums per chunk.
HTTP-level compression (\texttt{Accept-Encoding: gzip,deflate}) is
applied on top of either; clients negotiate this with the server in
the normal way.

\section{Subsetting}

Subsetting is the single most important capability OPeNDAP provides
\texttt{grep-dap}: a client can pull a few representative slices of a
remote variable at low cost, without downloading the full archive.
Three kinds of subset matter:

\paragraph{Variable subsetting.} List only the variables you want by
name. Under DAP4, list them by FQN. This alone can reduce a 1\,GB
file to a 1\,kB response if you only want one column, and~--- as the
pydap docs emphasize~--- it is much cheaper than fetching everything
and dropping variables client-side, because the server returns a
\emph{constrained DMR} containing only the requested variables.

\paragraph{Hyperslab subsetting (gridded data).} For each multi-dim
variable, append \texttt{[start:stride:stop]} per axis. Strided reads
work; non-contiguous reads do not (you must issue separate requests).
For shared-dimension datasets in DAP4, prefer the shared-dimension
form above.

\paragraph{Selection / filtering (sequence data).} For tabular data
in DAP2 \texttt{Sequence}s, expressions like \texttt{seq.depth\&seq.temp\&seq.profile\_id=5}
project two columns and filter rows. DAP4 has analogous semantics for
its sequence-like types but the conventions are still in flux; check
the server.

\paragraph{What about repeated holes?} OPeNDAP does not have a native
``random list of indices'' subset operator. The two ways forward are
(a) repeated single-element CEs, or (b) a bounding hyperslab plus
client-side masking; \texttt{pydap}/\texttt{xarray} use the latter.

\section{pydap: the Python client}

\texttt{pydap} is a pure-Python DAP client (and, separately, a server,
which is not the concern of \texttt{grep-dap}). As of version 3.5+ it
uses Python \texttt{requests} for transport, supports
\texttt{requests\_cache} for caching, covers most of the DAP2 and DAP4
data models, and is being developed by members of the OPeNDAP team
with coordination with Unidata.

\subsection{Opening a dataset}

\begin{lstlisting}[language=Python]
from pydap.client import open_url

# DAP2 (current default; subject to change)
ds = open_url("http://test.opendap.org/opendap/data/nc/coads_climatology.nc")

# DAP4 (recommended for modern servers)
ds = open_url("http://test.opendap.org/opendap/data/nc/coads_climatology.nc",
              protocol="dap4")

# Caching session (avoids re-fetching the DMR)
from pydap.net import create_session
session = create_session(use_cache=True)
ds = open_url(url, protocol="dap4", session=session)
\end{lstlisting}

\texttt{open\_url} fetches the metadata document (DDS+DAS for DAP2,
DMR for DAP4) and constructs a \texttt{DatasetType}~--- a dictionary-
like object containing lazy proxies. \emph{No data is fetched until a
variable is sliced.}

\subsection{Inspecting metadata}

\begin{lstlisting}[language=Python]
list(ds.keys())                  # top-level variable / group names
ds.tree()                        # full hierarchical view (DAP4)
ds['SST'].shape                  # shape, no data fetch
ds['SST'].dtype                  # numpy dtype
ds['SST'].attributes             # dict of attributes (DAS / DMR-attrs)
ds['SST'].units                  # lazy attribute access
ds.attributes                    # dataset-level (global) attributes
\end{lstlisting}

For DAP4 datasets with groups, FQNs apply:
\texttt{ds['/gt3r/heights/lat\_ph']}.

\subsection{Retrieving data}

Slicing a variable triggers a constrained request to the server:
\begin{lstlisting}[language=Python]
sst = ds['SST'][0, 50:60, 100:110]   # 1 time step, 10x10 lat-lon window
data = sst.data                       # numpy array
\end{lstlisting}
Under the hood, \texttt{pydap} builds the appropriate CE
(DAP2 hyperslab or DAP4 \texttt{dap4.ce}), URL-escapes it, performs
the GET, parses the binary response, and yields a numpy array.

\subsection{Subsetting through pydap}

Two equivalent levers:
\begin{itemize}
  \item \textbf{Python slicing.} As above. The user thinks in
    Python; \texttt{pydap} writes the CE. Easy and safe for most
    cases.
  \item \textbf{Explicit CE.} Pass \texttt{output\_grid=False} or
    build the constrained URL by hand and pass it to
    \texttt{open\_url}. Useful when you already know the FQNs and
    indices you want, or when \texttt{xarray}-style aggregation
    needs precise control.
\end{itemize}

For collections of files, \texttt{pydap} provides
\texttt{consolidate\_metadata(urls, concat\_dim=...)} which caches the
shared dimension arrays in SQLite for the session, reportedly
yielding 10--100$\times$ speedups when iterating over many files.

\subsection{pydap as an xarray backend}

\texttt{xarray.open\_dataset(url, engine="pydap")} (note the URL
scheme stays HTTP; do not use \texttt{dap4://}) makes the remote
dataset look like a local \texttt{xarray.Dataset}. \emph{Caveat from
the pydap docs:} \texttt{xarray}'s
\texttt{open\_mfdataset} does not, by default, push slices through to
the server for each file; it requests the chunk and slices locally.
To force per-file server-side subsetting, chunk the dataset so the
chunk size matches the desired output extent, or aggregate URLs
yourself. Overzealous chunking pushes the dask graph to thousands of
tiny chunks and can be 10$\times$ \emph{slower} than fetching more
data with fewer chunks.

\subsection{Authentication}

NASA Earthdata and many other archives require authentication. The
common patterns are:
\begin{itemize}
  \item A \texttt{.netrc} file with credentials for
    \texttt{urs.earthdata.nasa.gov}, picked up automatically.
  \item An \texttt{earthaccess}-managed session passed in via
    \texttt{session=}.
  \item A \texttt{requests.Session} carrying cookies or a bearer
    token.
\end{itemize}
For open archives (test.opendap.org, many academic THREDDS servers)
no authentication is needed.

\section{Implications for grep-dap}

A few things the protocol survey makes clear for the project:
\begin{enumerate}
  \item \textbf{Prefer DAP4 when available.} A single \texttt{.dmr}
    fetch yields both structure and attributes; FQNs eliminate
    name-collision ambiguity; chunked transfer and checksums make
    long fetches more robust. Fall back to DAP2 only when the server
    refuses DAP4.
  \item \textbf{Treat the DMR/DDS+DAS as the cheap probe.} Before
    pulling any bytes, \texttt{grep-dap} should fetch the metadata
    document and reason about variable names, shapes, dtypes,
    dimensions, groups, and whatever attributes are present. Most
    discriminating signal is here.
  \item \textbf{Use small samples deliberately.} For each candidate
    variable, a tiny hyperslab (say \texttt{[0:1:N]} along each axis
    with \texttt{N} small) is cheap; statistical fingerprints
    computed from such samples are how the contrastive-learning
    approach hooked into OPeNDAP in the first place.
  \item \textbf{Variable subsetting before data retrieval.} On
    multi-variable datasets (think L2 swaths with hundreds of
    variables), fetching the constrained DMR with only the
    variables of interest is dramatically cheaper than fetching the
    full DMR.
  \item \textbf{Cache the metadata aggressively.} Hitting the same
    server with the same DMR request repeatedly is wasteful and
    impolite; \texttt{pydap}'s session cache or
    \texttt{consolidate\_metadata} should be on by default.
  \item \textbf{Catalogs and directory structure are part of the
    signal.} Server catalog pages (Hyrax's \texttt{contents.html},
    THREDDS catalog XML) name datasets and group them; filenames
    and paths often encode platform, sensor, level, and date in
    informative ways that the loose semantic metadata may omit.
\end{enumerate}

\section{Open questions and caveats}

\begin{itemize}
  \item \texttt{pydap}'s documentation notes that several advanced
    sections (server-side functions, proxy configuration, some
    response modules) are explicitly flagged as ``may be outdated.''
    Anything in those sections should be re-verified before
    relying on it.
  \item \texttt{pydap} does not currently expose a generic
    ``list functions this server supports'' call; server-side
    functions are server-specific knowledge.
  \item \texttt{pydap} does not (yet) aggregate URLs the way
    \texttt{xarray} does; aggregation across files is handled by
    \texttt{xarray} with \texttt{pydap} as the per-URL backend.
  \item Some servers serve DAP4 but their DMRs decline to expose
    certain attributes that DAP2 would have provided (or vice
    versa). When metadata seems implausibly thin, try the other
    protocol on the same URL.
\end{itemize}

\section{References}

\begin{itemize}
  \item OPeNDAP, Inc.: \url{https://www.opendap.org/}
  \item DAP4 specification (Unidata/OPeNDAP, 2016):
    \url{https://opendap.github.io/dap4-specification/DAP4.html}
  \item DAP2 description, Cornillon et al.\ 1993:
    \url{https://zenodo.org/records/10610992}
  \item DAP2 protocol, Gallagher \& Flierl 1996:
    \url{https://zenodo.org/records/10654616}
  \item \texttt{pydap} documentation:
    \url{https://pydap.github.io/pydap/en/intro.html}
  \item \texttt{pydap} source / docs:
    \url{https://github.com/pydap/pydap}
  \item Hyrax data server:
    \url{https://www.opendap.org/software/hyrax-data-server/}
  \item THREDDS Data Server:
    \url{https://www.unidata.ucar.edu/software/tds/}
  \item NASA Earthdata OPeNDAP portal:
    \url{https://www.earthdata.nasa.gov/engage/open-data-services-software/earthdata-developer-portal/opendap}
  \item Test server (Hyrax, public):
    \url{http://test.opendap.org/opendap/}
\end{itemize}

\end{document}
